\documentclass[11pt,a4paper]{article}
% ---- Préambule commun ----
\usepackage[utf8]{inputenc}
\usepackage[T1]{fontenc}
\usepackage{lmodern}
\usepackage[french]{babel}
\usepackage[a4paper,margin=2.2cm,top=2.6cm,bottom=2.4cm]{geometry}
\usepackage{amsmath,amssymb,amsthm}
\usepackage{enumitem}
\usepackage{xcolor}
\usepackage{listings}
\usepackage{fancyhdr}
\usepackage[most]{tcolorbox}
\usepackage{hyperref}

\definecolor{teal}{HTML}{026573}
\definecolor{tealclair}{HTML}{E6F2F3}
\definecolor{grisclair}{HTML}{F5F5F5}
\definecolor{orangefonce}{HTML}{B35900}

\hypersetup{colorlinks=true,linkcolor=teal,urlcolor=teal,citecolor=teal}

\setlength{\parindent}{0pt}
\setlength{\parskip}{3pt}

% ---- Style Python ----
\lstdefinestyle{python}{
  language=Python,
  basicstyle=\ttfamily\small,
  keywordstyle=\color{teal}\bfseries,
  commentstyle=\color{gray}\itshape,
  stringstyle=\color{orangefonce},
  showstringspaces=false,
  numbers=left,
  numberstyle=\tiny\color{gray},
  numbersep=8pt,
  frame=leftline,
  framerule=1.2pt,
  rulecolor=\color{teal},
  xleftmargin=14pt,
  backgroundcolor=\color{grisclair},
  breaklines=true,
  literate={é}{{\'e}}1 {è}{{\`e}}1 {ê}{{\^e}}1 {à}{{\`a}}1 {ô}{{\^o}}1 {û}{{\^u}}1 {î}{{\^i}}1 {ç}{{\c c}}1
}
\lstset{style=python}

% ---- Compteur d'exercices ----
\newcounter{exo}
\newcommand{\exo}[1]{%
  \stepcounter{exo}%
  \vspace{10pt}%
  \noindent\begin{tcolorbox}[colback=tealclair,colframe=teal,boxrule=1pt,
      left=8pt,right=8pt,top=5pt,bottom=5pt,arc=2pt]
    \textbf{\large Exercice \theexo} \;---\; \textbf{#1}
  \end{tcolorbox}%
  \vspace{4pt}%
  \nopagebreak
}

\newcommand{\partiecours}[1]{%
  \vspace{14pt}%
  {\color{teal}\hrule height 2pt}%
  \vspace{4pt}%
  {\Large\bfseries\color{teal} #1}%
  \vspace{2pt}%
  {\color{teal}\hrule height 0.6pt}%
  \vspace{8pt}%
}

\newtcolorbox{rappel}[1][Rappel]{
  colback=white,colframe=gray!60,boxrule=0.8pt,
  fonttitle=\bfseries,title=#1,coltitle=black,colbacktitle=grisclair,
  left=6pt,right=6pt}

\newtcolorbox{indication}{
  colback=white,colframe=orangefonce!60,boxrule=0.8pt,
  left=6pt,right=6pt,top=4pt,bottom=4pt}

\newcommand{\reponse}{\vspace{4pt}\noindent{\color{teal}\textbf{Solution.}}\;}

\setlength{\headheight}{14pt}
\pagestyle{fancy}
\fancyhf{}
\fancyfoot[C]{\thepage}
\renewcommand{\headrulewidth}{0.4pt}
\renewcommand{\footrulewidth}{0pt}

\fancyhead[L]{\small\color{teal} TD --- Complexité algorithmique}
\fancyhead[R]{\small\color{teal} Corrigé}

\begin{document}

\begin{center}
  {\color{teal}\hrule height 2.5pt}
  \vspace{10pt}
  {\Huge\bfseries\color{teal} Corrigé du TD}\\[6pt]
  {\LARGE\bfseries Calcul de complexité algorithmique}\\[10pt]
  {\large Cas itératif \;$\bullet$\; Cas récursif}\\[8pt]
  {\normalsize CPGE --- filières MP / PSI \quad$|$\quad Informatique pour tous}\\[6pt]
  \vspace{6pt}
  {\color{teal}\hrule height 1pt}
\end{center}

\vspace{8pt}

% =====================================================================
\partiecours{Partie I --- Le cas itératif : boucles \texttt{for} imbriquées}
% =====================================================================

\exo{Échauffement --- compter avant de majorer}
\reponse
\begin{enumerate}[label=(\alph*),leftmargin=*]
  \item $n$ exécutions. Complexité $\Theta(n)$.
  \item $\displaystyle\sum_{i=0}^{n-1}\sum_{j=0}^{n-1}1 = n^2$. Complexité $\Theta(n^2)$.
  \item $m\,n$ exécutions. Complexité $\Theta(mn)$.
  \item $\displaystyle\sum_{i=0}^{n-1}\sum_{j=0}^{4}\sum_{k=0}^{n-1}1 = 5n^2$.
        Complexité $\Theta(n^2)$ : la boucle de longueur fixe $5$ n'apporte qu'une
        constante multiplicative, invisible dans le $\Theta$.
\end{enumerate}

\textbf{Synthèse.} Non : $mn$ n'est $\Theta(n^2)$ que si $m=\Theta(n)$.
Si $m$ est une constante, la complexité est $\Theta(n)$ ; si $m=n^2$, elle est
$\Theta(n^3)$. Dès qu'un algorithme dépend de \emph{plusieurs} paramètres, la complexité
doit s'exprimer en fonction de \emph{tous} ces paramètres : écrire $\Theta(mn)$ et non
$\Theta(n^2)$.

% ---------------------------------------------------------------------
\exo{Boucles triangulaires}
\reponse
\textbf{1.} Pour (a) : $\displaystyle\sum_{i=0}^{n-1} i = \frac{n(n-1)}{2}$.
Pour (b) : la boucle interne \texttt{range(i,n)} fait $n-i$ tours, d'où
$\displaystyle\sum_{i=0}^{n-1}(n-i) = \sum_{j=1}^{n} j = \frac{n(n+1)}{2}$.
La différence vaut $\frac{n(n+1)}{2}-\frac{n(n-1)}{2}= n$ : c'est exactement le cas
$j=i$, présent dans (b) et absent de (a). Dans les deux cas la complexité est
$\Theta(n^2)$, avec la même constante dominante $\tfrac12$.

\textbf{2.} La boucle exécute \texttt{s += 1} une fois par triplet $(i,j,k)$ tel que
$1\le i\le j\le k\le n$. Le nombre de tels triplets est le nombre de combinaisons
\emph{avec répétition} de $3$ éléments parmi $n$ :
\[
  \binom{n+3-1}{3} = \binom{n+2}{3} = \frac{n(n+1)(n+2)}{6}\;\sim\;\frac{n^3}{6}.
\]
Complexité $\Theta(n^3)$.

\smallskip
\emph{Vérification directe.} $\displaystyle\sum_{i=1}^{n}\sum_{j=i}^{n}(n-j+1)
=\sum_{i=1}^{n}\sum_{\ell=1}^{n-i+1}\ell
=\sum_{i=1}^{n}\frac{(n-i+1)(n-i+2)}{2}
=\sum_{p=1}^{n}\frac{p(p+1)}{2}=\frac{n(n+1)(n+2)}{6}.$

\textbf{3.} Pour $p$ boucles imbriquées avec $1\le i_1\le\dots\le i_p\le n$, le nombre
d'itérations est $\displaystyle\binom{n+p-1}{p}\sim\frac{n^p}{p!}$, soit $\Theta(n^p)$
à $p$ fixé.

% ---------------------------------------------------------------------
\exo{Indices strictement croissants}
\reponse
\textbf{1.} Dans le pire des cas (aucun triplet de somme nulle), le test est effectué
une fois par triplet $0\le i<j<k\le n-1$, soit
\[
  \binom{n}{3} = \frac{n(n-1)(n-2)}{6}\sim\frac{n^3}{6}\quad\text{tests},
  \qquad\text{complexité } \Theta(n^3).
\]

\textbf{2.} Dans le meilleur des cas, \texttt{T[0]+T[1]+T[2] == 0} et l'algorithme
renvoie après un seul test : complexité $\Theta(1)$.
Les complexités du meilleur et du pire cas n'étant pas du même ordre, il n'existe
\emph{aucune} fonction $g$ telle que le coût soit $\Theta(g(n))$ pour toute entrée.
On doit donc préciser : « $\Theta(n^3)$ \emph{dans le pire des cas} », ou se contenter
de « $O(n^3)$ ».

\textbf{3.} Pour (b) : $\displaystyle\sum_{i=0}^{n-1}\sum_{j=i+1}^{n-1} n
= n\binom{n}{2} = \frac{n^2(n-1)}{2}\sim\frac{n^3}{2}$, soit $\Theta(n^3)$.
L'ordre de grandeur est le même que (a) ; seule la constante change ($\tfrac12$ au lieu
de $\tfrac16$). \emph{Conclusion utile :} remplacer \texttt{range(j+1,n)} par
\texttt{range(n)} ne fait pas changer de classe de complexité.

% ---------------------------------------------------------------------
\exo{Bornes non linéaires et pas variable}
\reponse
\textbf{(a)} $\displaystyle\sum_{i=0}^{n-1} i^2 = \frac{(n-1)n(2n-1)}{6}\sim\frac{n^3}{3}$.
Complexité $\Theta(n^3)$.

\textbf{(b)} La boucle \texttt{range(1,n+1,i)} parcourt $1, 1+i, 1+2i,\dots$ et effectue
$\big\lceil n/i\big\rceil$ tours. Donc
\[
  N(n)=\sum_{i=1}^{n}\Big\lceil \frac n i\Big\rceil,
  \qquad \sum_{i=1}^{n}\frac n i \;\le\; N(n)\;\le\; \sum_{i=1}^{n}\Big(\frac n i +1\Big)
  = n H_n + n .
\]
Comme $H_n = \ln n + \gamma + o(1)$, on obtient $N(n) \sim n\ln n$, soit
$\boxed{\Theta(n\log n)}$.

\textbf{(c)} La boucle \texttt{while} interne s'exécute tant que $j^2\le i$ ; elle fait
donc $\#\{j\ge 0 : j^2\le i\} = \lfloor\sqrt i\rfloor + 1$ tours. D'où
\[
  N(n)=\sum_{i=1}^{n}\big(\lfloor\sqrt i\rfloor+1\big).
\]
Par comparaison série–intégrale, $\displaystyle\int_0^{n}\!\sqrt t\,dt \le
\sum_{i=1}^n \sqrt i \le \int_1^{n+1}\!\sqrt t\,dt$, donc
$\sum_{i=1}^n\sqrt i \sim \frac23 n^{3/2}$ et
$N(n)\sim \frac23 n^{3/2}$ : complexité $\Theta\!\left(n^{3/2}\right)$.

% ---------------------------------------------------------------------
\exo{Progression géométrique : attention aux fausses évidences}
\reponse
\textbf{1.} Avec $n=2^p$, la variable \texttt{i} prend les valeurs
$1,2,4,\dots,2^{p-1}$ (elle s'arrête dès que $i\ge n$). Le coût total est
\[
  \sum_{t=0}^{p-1} 2^{t} = 2^{p}-1 = n-1 .
\]
Complexité $\Theta(n)$, et non $\Theta(n\log n)$ : la boucle externe fait bien
$\log_2 n$ tours, mais les tours ont des coûts très inégaux et la somme géométrique est
dominée par son \emph{dernier} terme. Le résultat n'est donc pas surprenant une fois
la somme écrite --- il l'est seulement si l'on multiplie hâtivement.

\textbf{2.} Ici chaque tour coûte $n$, indépendamment de \texttt{i} :
coût total $= p\cdot n = n\log_2 n$, soit $\Theta(n\log n)$.

\textbf{3.} \emph{Moralité.} Le coût d'une boucle est la \textbf{somme} des coûts de ses
tours, jamais le produit « nombre de tours $\times$ coût du dernier tour ». Ce produit
n'est une bonne approximation que lorsque les tours ont un coût du même ordre (cas (b)) ;
il surestime lourdement dans le cas d'une progression géométrique (cas (a)).

\textbf{4.} Avec \texttt{i = i*i} et $i_0=2$ : après $t$ tours $i = 2^{2^{t}}$.
La boucle s'arrête au premier $t=K$ tel que $2^{2^{K}}\ge n$, donc elle effectue
$\Theta(\log\log n)$ tours. Le coût total vaut
$\displaystyle\sum_{t=0}^{K-1} 2^{2^{t}}$, somme à croissance doublement exponentielle,
donc dominée par son dernier terme :
\[
  2^{2^{K-1}} \;\le\; \sum_{t=0}^{K-1} 2^{2^{t}} \;\le\; 2\cdot 2^{2^{K-1}} .
\]
Or $2^{2^{K-1}} < n$ et $\big(2^{2^{K-1}}\big)^2 = 2^{2^{K}}\ge n$, donc
$\sqrt n \le 2^{2^{K-1}} < n$. Le coût total est donc $O(n)$ et $\Omega(\sqrt n)$,
mais il n'existe pas de $c$ tel que ce coût soit $\Theta(n^{c})$ : il oscille selon la
position de $n$ entre deux valeurs consécutives $2^{2^{t}}$.

% =====================================================================
\partiecours{Partie II --- Le cas itératif : boucles \texttt{while}}
% =====================================================================

\exo{Les trois schémas fondamentaux}
\reponse
\textbf{1.}
\begin{itemize}[leftmargin=*,itemsep=2pt]
  \item \textbf{(a)} après $k$ tours $x = n-k$ ; arrêt quand $n-k=0$ :
        exactement $n$ tours, $\Theta(n)$.
  \item \textbf{(b)} après $k$ tours $x=n-kc$ ; arrêt au premier $k$ tel que $n-kc\le0$,
        soit $k=\lceil n/c\rceil$ tours, $\Theta(n)$.
  \item \textbf{(c)} on montre par récurrence que $x = \lfloor n/2^{k}\rfloor$ après
        $k$ tours (on utilise $\big\lfloor\lfloor n/2\rfloor/2\big\rfloor=\lfloor n/4\rfloor$).
        L'arrêt a lieu au premier $k$ tel que $\lfloor n/2^{k}\rfloor \le 1$, c'est-à-dire
        $n/2^{k}<2$, soit $2^{k}>n/2$, soit $k>\log_2 n-1$ :
        le nombre de tours est $\lfloor \log_2 n\rfloor$. Complexité $\Theta(\log n)$.
  \item \textbf{(d)} après $k$ tours $x=3^{k}$ ; arrêt quand $3^{k}\ge n$ :
        $k=\lceil\log_3 n\rceil$ tours, $\Theta(\log n)$ (la base du logarithme est une
        constante multiplicative : $\log_3 n = \log_2 n/\log_2 3$).
\end{itemize}

\textbf{2.} Le nombre de tours $\lceil n/c\rceil$ dépend de $c$, mais $1/c$ est une
constante multiplicative : pour $c$ \emph{fixé}, la complexité reste $\Theta(n)$.
Diviser $c$ par deux ne change pas la classe de complexité, seulement la constante.

\textbf{3.} Si $t=\lfloor\log_2 n\rfloor+1$ est le nombre de bits de $n$, alors
$n\approx 2^{t}$ et :
\begin{itemize}[leftmargin=*,itemsep=1pt]
  \item (c) coûte $\Theta(\log n) = \Theta(t)$ : \emph{linéaire} en la taille de l'entrée ;
  \item (a) coûte $\Theta(n) = \Theta(2^{t})$ : \emph{exponentiel} en la taille de
        l'entrée.
\end{itemize}
C'est la distinction fondamentale entre complexité \emph{pseudo-polynomiale} (polynomiale
en la \emph{valeur} de l'entrée) et complexité polynomiale (en la \emph{taille} de
l'entrée). L'algorithme (a) est en pratique inutilisable pour un entier de $200$ bits.

% ---------------------------------------------------------------------
\exo{Algorithme d'Euclide}
\reponse
\textbf{1.} Après un tour, le couple $(a,b)$ devient $(b,\,a\bmod b)$. Par définition du
reste, $0\le a\bmod b< b$. Le \emph{variant} $V = b$ est donc un entier naturel
strictement décroissant à chaque tour : la boucle termine.

\textbf{2.} Soit $a>b>0$.
\begin{itemize}[leftmargin=*,itemsep=1pt]
  \item Si $b\le a/2$ : $a\bmod b < b \le a/2$.
  \item Si $b>a/2$ : alors $a/b<2$ et $a>b$, donc le quotient de $a$ par $b$ vaut
        exactement $1$ et $a\bmod b = a-b < a - a/2 = a/2$.
\end{itemize}
Dans les deux cas $a \bmod b < a/2$.

\textbf{3.} Notons $(a_0,b_0)=(a,b)$ et $(a_{k+1},b_{k+1})=(b_k,\,a_k\bmod b_k)$.
Alors $a_{k+2}=b_{k+1}=a_k \bmod b_k < a_k/2$ d'après \textbf{2}.
Donc $a_{2j} < a_0/2^{j}$. Comme $a_{2j}\ge 1$ tant que la boucle tourne, on a
$2^{j}< a_0$, soit $j<\log_2 a$ : le nombre de tours est $\le 2\log_2 a$.
En remarquant qu'après le premier tour le premier argument vaut $b=\min(a,b)$ (si
$a>b$), on obtient bien $O(\log \min(a,b))$.

\textbf{4.} Pour $k\ge2$, $F_{k+1}=F_k+F_{k-1}$ avec $0\le F_{k-1}<F_k$, donc
$F_{k+1}\bmod F_k = F_{k-1}$. L'appel \texttt{pgcd}$(F_{k+1},F_k)$ produit donc la suite
de couples
\[
  (F_{k+1},F_k)\to(F_k,F_{k-1})\to\cdots\to(F_2,F_1)=(1,1)\to(1,0),
\]
soit exactement $k$ tours de boucle. La borne logarithmique est donc atteinte à une
constante près. Comme $F_k \sim \varphi^{k}/\sqrt5$, on a
$k \sim \log_\varphi\!\big(\sqrt5\,F_k\big) = \dfrac{\ln a}{\ln\varphi}
\approx 1{,}44\,\log_2 a$ : c'est le \emph{théorème de Lamé}. Le pire des cas d'Euclide
est atteint sur deux nombres de Fibonacci consécutifs.

% ---------------------------------------------------------------------
\exo{Exponentiation rapide}
\reponse
\textbf{1.} \emph{Invariant} : à chaque passage au test de boucle, $r\cdot x^{\,n}
= x_0^{\,n_0}$.
\emph{Initialisation} : $r=1$, $x=x_0$, $n=n_0$ : vrai.
\emph{Conservation} : si $n$ est pair, $(r,x,n)\mapsto (r,\,x^2,\,n/2)$ et
$r\,(x^2)^{n/2}=r\,x^{n}$. Si $n$ est impair, $(r,x,n)\mapsto(rx,\,x^2,\,(n-1)/2)$ et
$rx\,(x^2)^{(n-1)/2}= r\,x^{n}$. L'invariant est conservé.
\emph{Terminaison} : $n$ est un entier naturel strictement décroissant (pour $n\ge1$,
$\lfloor n/2\rfloor<n$). À la sortie $n=0$, donc $r = r\cdot x^{0} = x_0^{\,n_0}$.

\textbf{2.} Après $k$ tours, $n$ vaut $\lfloor n_0/2^{k}\rfloor$ ; la boucle s'arrête
au premier $k$ tel que $\lfloor n_0/2^{k}\rfloor = 0$, soit $2^{k}>n_0$ :
le nombre de tours est $\lfloor\log_2 n_0\rfloor+1$, c'est-à-dire le nombre de bits
de $n_0$.

\textbf{3.} Chaque tour effectue exactement une élévation au carré \texttt{x = x*x},
soit $\lfloor\log_2 n\rfloor+1$ multiplications. La multiplication \texttt{r = r*x}
n'est effectuée que lorsque le bit courant de $n$ vaut $1$, soit $\nu(n)$ fois au total
(les bits successifs de $n$ sont exactement les valeurs de \texttt{n \% 2}). D'où
\[
  M(n) = \lfloor\log_2 n\rfloor + 1 + \nu(n).
\]
\emph{Vérification} : $n=1$ donne $0+1+1=2$ (une multiplication dans \texttt{r}, un
carré) ; $n=3$ donne $1+1+2=4$.
\emph{Économie possible} : au dernier tour, l'élévation \texttt{x = x*x} n'est jamais
réutilisée ; on peut donc la supprimer et gagner une multiplication, d'où
$M'(n)=\lfloor\log_2 n\rfloor+\nu(n)$.

\textbf{4.} Comme $1\le\nu(n)\le\lfloor\log_2 n\rfloor+1$,
\[
  \log_2 n + 2 \;\le\; M(n) \;\le\; 2\big(\lfloor\log_2 n\rfloor+1\big).
\]
La complexité est donc $\Theta(\log n)$ \emph{dans tous les cas}.
Le pire cas ($\nu$ maximal) est atteint pour $n=2^{p}-1$ (écriture $11\cdots1$),
le meilleur ($\nu(n)=1$) pour $n=2^{p}$. À comparer aux $n-1$ multiplications de
l'algorithme naïf.

% ---------------------------------------------------------------------
\exo{Deux indices qui se croisent}
\reponse
\textbf{1.} \emph{Variant} : $V = d-g$. C'est un entier, positif tant que la condition
$g<d$ est vraie, et chaque tour effectue soit \texttt{g += 1} soit \texttt{d -= 1}
(ou renvoie), donc $V$ décroît strictement de $1$. La boucle termine.

\textbf{2.} Initialement $V=n-1$ ; $V$ décroît de $1$ par tour et la boucle s'arrête dès
que $V=0$. Il y a donc au plus $n-1$ tours. La borne est atteinte lorsque
l'algorithme renvoie \texttt{None} (aucune paire de somme $s$), par exemple
$T=[0,1,2,\dots,n-1]$ et $s$ strictement supérieur à $T[n-2]+T[n-1]$.

\textbf{3.} Chaque tour coûte $\Theta(1)$ : la complexité dans le pire des cas est
$\Theta(n)$. L'algorithme naïf à deux boucles \texttt{for} imbriquées teste les
$\binom n2$ paires, soit $\Theta(n^2)$ : le tri du tableau (supposé acquis ici) permet
de passer du quadratique au linéaire.

\textbf{4.} \emph{Invariant} : « si une solution $(i,j)$ avec $i<j$ et $T[i]+T[j]=s$
existe, alors $g\le i$ et $j\le d$ ».
\emph{Initialisation} : $g=0$ et $d=n-1$ encadrent tout indice.
\emph{Conservation} : supposons l'invariant vrai et le tour non concluant.
\begin{itemize}[leftmargin=*,itemsep=1pt]
  \item Si $T[g]+T[d]<s$ : montrons $g<i$. Par l'absurde, si $g=i$, alors
        $T[i]+T[d]<s=T[i]+T[j]$ donne $T[d]<T[j]$, donc $d<j$ (tableau trié
        strictement croissant), ce qui contredit $j\le d$. Donc $g<i$, et après
        \texttt{g += 1} on a encore $g\le i$ ; $d$ est inchangé.
  \item Si $T[g]+T[d]>s$ : symétriquement $d>j$, donc après \texttt{d -= 1} on a
        encore $j\le d$.
\end{itemize}
\emph{Conclusion} : si l'algorithme renvoie \texttt{None}, c'est que $g=d$ alors qu'une
solution imposerait $g\le i<j\le d$, donc $g<d$ : contradiction. Il n'existe donc pas de
solution. L'algorithme est correct.

% =====================================================================
\partiecours{Partie III --- Le cas récursif : équation caractéristique}
% =====================================================================

\exo{Tours de Hanoï et variantes}
\reponse
\textbf{1.} $T(n) = 2\,T(n-1) + 1$ pour $n\ge1$, et $T(0)=0$.

\textbf{2.} \emph{Homogène} : $r = 2$, solution $A\,2^{n}$.
\emph{Particulière} : le second membre est constant et $1$ n'est pas racine de
$r=2$ ; on cherche $T_p = c$ : $c = 2c+1 \Rightarrow c=-1$.
Donc $T(n)=A\,2^{n}-1$, et $T(0)=0$ donne $A=1$ :
\[
  \boxed{T(n) = 2^{n}-1},\qquad T(n)=\Theta(2^{n}).
\]
Vérification : $T(1)=1$, $T(2)=3$, $T(3)=7$. \checkmark

\textbf{3.}
\begin{enumerate}[label=(\alph*),leftmargin=*,itemsep=3pt]
  \item $T(n)=T(n-1)+n$ : par télescopage $T(n)=\sum_{k=1}^{n}k=\dfrac{n(n+1)}{2}$,
        soit $\Theta(n^2)$.
        \emph{(Méthode caractéristique : racine $1$ simple, second membre polynomial de
        degré $1$ ; on cherche $T_p$ de degré $1+1=2$.)}
  \item $T(n)=T(n-1)+n^2$ : $T(n)=\displaystyle\sum_{k=1}^{n}k^2=\frac{n(n+1)(2n+1)}{6}$,
        soit $\Theta(n^3)$.
  \item $T(n)=2T(n-1)+n$ : homogène $A2^{n}$ ; particulière $\alpha n+\beta$ donne
        $\alpha n + \beta = 2\alpha(n-1)+2\beta+n$, d'où $\alpha=-1$ et $\beta=-2$.
        Avec $T(0)=0$ : $A=2$ et
        $\boxed{T(n)=2^{n+1}-n-2}$, soit $\Theta(2^{n})$.
        \emph{(Vérification : $T(1)=4-3=1=2\cdot0+1$. \checkmark)}
  \item $T(n)=3T(n-1)-2T(n-2)$ : équation caractéristique $r^2-3r+2=0$, racines
        $1$ et $2$. Donc $T(n)=A+B\,2^{n}$ ; $T(0)=1$ et $T(1)=3$ donnent $A=-1$, $B=2$ :
        $\boxed{T(n)=2^{n+1}-1}$, soit $\Theta(2^{n})$.
\end{enumerate}

% ---------------------------------------------------------------------
\exo{Coût de Fibonacci naïf}
\reponse
\textbf{1.} Un appel \texttt{fib(n)} avec $n\ge2$ engendre lui-même plus les appels
issus de \texttt{fib(n-1)} et \texttt{fib(n-2)} :
\[
  C(n) = C(n-1)+C(n-2)+1\quad (n\ge2),\qquad C(0)=C(1)=1.
\]

\textbf{2.} \emph{Homogène} : $r^2 = r+1$, racines $\varphi=\frac{1+\sqrt5}{2}$ et
$\psi=\frac{1-\sqrt5}{2}$ (noter $\varphi\psi=-1$, $\varphi+\psi=1$).
\emph{Particulière} : $1$ n'est pas racine ($1\ne 1+1$), on cherche $c$ constante :
$c=2c+1 \Rightarrow c=-1$. Donc
\[
  C(n) = A\varphi^{\,n}+B\psi^{\,n}-1 .
\]
Les conditions $C(0)=C(1)=1$ donnent $A+B=2$ et $A\varphi+B\psi=2$, d'où
$A = \dfrac{2\varphi}{\sqrt5}$ et $B = -\dfrac{2\psi}{\sqrt5}$.

\textbf{3.} Ainsi
\[
  C(n) = \frac{2}{\sqrt5}\big(\varphi^{\,n+1}-\psi^{\,n+1}\big)-1 = 2F_{n+1}-1,
\]
puisque $F_m = \dfrac{\varphi^{m}-\psi^{m}}{\sqrt5}$ (formule de Binet).
Vérification : $C(0)=2F_1-1=1$, $C(2)=2F_3-1=3$. \checkmark

\textbf{4.} Comme $|\psi|<1$, $\psi^{\,n+1}\to0$ et
$C(n)\sim \dfrac{2}{\sqrt5}\,\varphi^{\,n+1}$.
La complexité de \texttt{fib} est donc $\Theta(\varphi^{\,n})$ avec
$\varphi\approx1{,}618$ : elle est \emph{exponentielle}.

\textbf{5.} $C(80)=2F_{81}-1 = 75\,778\,124\,746\,287\,811 \approx 7{,}6\cdot10^{16}$
appels. À $10^{9}$ appels par seconde : $\approx 7{,}6\cdot10^{7}$ secondes,
soit environ \textbf{2,4 ans}. (À comparer aux $81$ additions de la version itérative.)

\textbf{6.} Avec mémoïsation, chaque valeur \texttt{fib(k)} pour $0\le k\le n$ n'est
calculée qu'une seule fois ; les appels suivants sont des lectures en $\Theta(1)$.
Le nombre d'appels « réellement calculés » est $n+1$ et chaque appel effectue un travail
$\Theta(1)$ (hors coût des grands entiers) : la complexité devient $\Theta(n)$, pour un
espace mémoire $\Theta(n)$.

% ---------------------------------------------------------------------
\exo{Racine double et second membre}
\reponse
\begin{enumerate}[label=(\alph*),leftmargin=*,itemsep=6pt]
  \item $r^2-4r+4=(r-2)^2$ : racine \emph{double} $r=2$, donc
        $T(n)=(A+Bn)2^{n}$. $T(0)=A=1$ ; $T(1)=(1+B)\cdot2=6\Rightarrow B=2$.
        \[\boxed{T(n)=(2n+1)\,2^{n}} ,\qquad T(n)=\Theta(n\,2^{n}).\]

  \item $r^2-2r+1=(r-1)^2$ : racine double $r=1$, homogène $A+Bn$.
        Le second membre étant constant et $1$ étant racine \emph{double}, on cherche
        $T_p=\alpha n^{2}$ :
        \[
          \alpha n^2 = 2\alpha(n-1)^2-\alpha(n-2)^2+2 = \alpha n^2 -2\alpha+2
          \;\Longrightarrow\; \alpha=1 .
        \]
        Donc $T(n)=A+Bn+n^2$ ; $T(0)=0\Rightarrow A=0$, $T(1)=B+1=1\Rightarrow B=0$ :
        \[\boxed{T(n)=n^{2}},\qquad \Theta(n^2).\]
        Vérification : $T(2)=2\cdot1-0+2=4$, $T(3)=2\cdot4-1+2=9$. \checkmark

  \item C'est exactement la récurrence de l'exercice~11 :
        $\boxed{T(n)=2F_{n+1}-1}$, soit $\Theta(\varphi^{\,n})$.

  \item $T(n)=2T(n-1)+2^{n}$ : l'homogène est $A\,2^{n}$ et $2$ \emph{est} racine simple ;
        on cherche donc $T_p=\alpha n 2^{n}$ :
        \[
          \alpha n 2^{n} = 2\alpha(n-1)2^{n-1}+2^{n} = \alpha(n-1)2^{n}+2^{n}
          \;\Longrightarrow\; \alpha = 1 .
        \]
        Donc $T(n)=(A+n)2^{n}$ et $T(0)=A=1$ :
        \[\boxed{T(n)=(n+1)\,2^{n}},\qquad \Theta(n\,2^{n}).\]
        Vérification : $T(1)=2\cdot1+2=4=2\cdot2^{1}$. \checkmark
\end{enumerate}

% =====================================================================
\partiecours{Partie IV --- Le cas récursif : récurrences de partition}
% =====================================================================

\exo{Déroulement à la main --- trois schémas de base}
\reponse
On pose $n=2^{k}$, donc $k=\log_2 n$.

\textbf{1. Recherche dichotomique.} $T(2^{k}) = T(2^{k-1})+1$, $T(1)=1$. En déroulant :
\[
  T(2^{k}) = T(2^{k-1})+1 = T(2^{k-2})+2 = \dots = T(2^{0}) + k = k+1 .
\]
Donc $T(n)=\log_2 n + 1 = \Theta(\log n)$.

\textbf{2. Somme par dichotomie.} L'appel sur un segment de taille $n$ effectue deux
appels sur des segments de taille $n/2$, plus un calcul de milieu et une addition :
$T(n)=2T(n/2)+c$ avec $c>0$ constante et $T(1)=1$. Déroulement :
\[
  T(2^{k}) = 2T(2^{k-1})+c = 4T(2^{k-2})+2c+c = \dots
  = 2^{k}T(1) + c\sum_{i=0}^{k-1}2^{i} = 2^{k} + c(2^{k}-1).
\]
Donc $T(n) = (1+c)n - c = \Theta(n)$.
\emph{Non}, cette version n'est pas meilleure que la boucle \texttt{for} naïve : même
classe $\Theta(n)$, mais avec une constante plus grande (gestion de la pile d'appels) et
un espace mémoire $\Theta(\log n)$ au lieu de $\Theta(1)$. Diviser pour régner n'apporte
rien lorsque la recombinaison ne réduit pas le travail.

\textbf{3.} $T(n)=2T(n/2)+n$, $T(1)=1$. Avec l'arbre des appels : au niveau $i$
($0\le i\le k$) il y a $2^{i}$ appels portant chacun sur une taille $n/2^{i}$, d'où un
coût de niveau
\[
  2^{i}\times \frac{n}{2^{i}} = n \quad\text{(identique à tous les niveaux)} .
\]
Il y a $k=\log_2 n$ niveaux internes plus le niveau des feuilles ($n$ feuilles de coût
$1$), d'où
\[
  T(n) = n\log_2 n + n = \Theta(n\log n).
\]

\textbf{4. Minimum et maximum simultanés.} On découpe en deux moitiés, on obtient
$(\min_1,\max_1)$ et $(\min_2,\max_2)$, puis $2$ comparaisons suffisent pour recombiner :
\[
  C(n)=2\,C(n/2)+2 \quad (n\ge4),\qquad C(2)=1 .
\]
Posons $n=2^{k}$ et $D(k)=C(2^{k})$ : $D(k)=2D(k-1)+2$, $D(1)=1$.
L'homogène est $A2^{k}$, la solution particulière constante vérifie $c=2c+2$, soit
$c=-2$. Donc $D(k)=A2^{k}-2$ et $D(1)=2A-2=1$ donne $A=\tfrac32$ :
\[
  D(k)= 3\cdot 2^{k-1}-2 \quad\Longrightarrow\quad \boxed{C(n)=\frac{3n}{2}-2}.
\]
L'algorithme itératif naïf (un parcours pour le maximum, un pour le minimum sur les
$n-1$ éléments restants) effectue $2n-3$ comparaisons. Le gain est d'environ $25\,\%$ :
le nombre de comparaisons reste $\Theta(n)$, seule la constante s'améliore --- ce qui est
justement l'objet de ce type d'analyse fine.

% ---------------------------------------------------------------------
\exo{Théorème général (Master Theorem)}
\reponse
\textbf{1.}
\begin{enumerate}[label=(\alph*),leftmargin=*,itemsep=4pt]
  \item $a=9$, $b=3$, $\alpha=\log_3 9=2$. $f(n)=n=O(n^{2-\varepsilon})$ avec
        $\varepsilon=1$ : \textbf{cas 1}, donc $T(n)=\Theta(n^{2})$.
  \item $a=1$, $b=3/2$, $\alpha=\log_{3/2}1=0$. $f(n)=1=\Theta(n^{0}\log^{0}n)$ :
        \textbf{cas 2} avec $c=0$, donc $T(n)=\Theta(\log n)$.
  \item $a=4$, $b=2$, $\alpha=2$. $f(n)=n^{2}=\Theta(n^{2}\log^{0}n)$ :
        \textbf{cas 2} avec $c=0$, donc $T(n)=\Theta(n^{2}\log n)$.
  \item $a=3$, $b=4$, $\alpha=\log_4 3\approx0{,}7925$.
        $f(n)=n\log n = \Omega(n^{0{,}7925+\varepsilon})$ pour $\varepsilon=0{,}1$ par
        exemple. Régularité : $3\,f(n/4)=\tfrac34\,n\log(n/4)\le\tfrac34\,n\log n
        = \tfrac34 f(n)$ avec $c=\tfrac34<1$ : \textbf{cas 3}, donc
        $T(n)=\Theta(n\log n)$.
  \item $a=2$, $b=2$, $\alpha=1$, $f(n)=n/\log n$. \textbf{Aucun cas ne s'applique}
        (voir \textbf{2}). Le résultat correct est $T(n)=\Theta(n\log\log n)$ :
        au niveau $i$ de l'arbre, $2^{i}$ appels de taille $n/2^{i}$ coûtent chacun
        $\dfrac{n/2^{i}}{\log n - i}$, soit un coût de niveau $\dfrac{n}{\log n - i}$ ;
        en sommant sur $i=0,\dots,\log n-1$ on obtient $n\,H_{\log n}\sim n\ln\log n$.
\end{enumerate}

\textbf{2.} Pour (e), $f(n)=n/\log n$ est \emph{plus petit} que $n^{\alpha}=n$, donc le
cas 3 est exclu ; mais il n'est pas \emph{polynomialement} plus petit : pour tout
$\varepsilon>0$, $\dfrac{n/\log n}{n^{1-\varepsilon}} = \dfrac{n^{\varepsilon}}{\log n}
\to+\infty$, donc $f\ne O(n^{1-\varepsilon})$ et le cas 1 est exclu. Enfin
$f(n)=n\log^{-1}n$ correspondrait au cas 2 avec $c=-1<0$, ce qui est interdit.
$f$ tombe donc dans la « zone grise » entre les cas 1 et 2 : le théorème est muet,
il faut revenir à l'arbre des appels.

\textbf{3.} $T(n)=2T(n/2)+n\log n$, $n=2^{k}$. Au niveau $i$ :
$2^{i}$ appels de taille $n/2^{i}$, de coût $\dfrac{n}{2^{i}}\log\dfrac{n}{2^{i}}$
chacun, soit un coût de niveau
\[
  2^{i}\cdot\frac{n}{2^{i}}\big(\log n - i\big) = n\,(k-i).
\]
En sommant sur $i=0,\dots,k-1$ :
\[
  T(n)=n\sum_{i=0}^{k-1}(k-i) = n\,\frac{k(k+1)}{2} = \Theta\!\big(n\log^{2}n\big).
\]
\emph{Remarque :} ce cas relève en fait du \textbf{cas 2} avec $c=1$
($f(n)=\Theta(n^{\alpha}\log n)$ avec $\alpha=1$), qui prédit bien
$\Theta(n\log^{2}n)$. L'arbre le confirme.

\textbf{4.} Le cas 3 affirme que le coût est dominé par la \emph{racine} de l'arbre.
Encore faut-il que le coût décroisse géométriquement quand on descend, ce qu'exprime
la condition $a\,f(n/b)\le c\,f(n)$ avec $c<1$. Avec $f(n)=n^{2}\sin^{2}n$ et $a=b=2$,
$f$ s'annule (ou devient minuscule) pour certaines valeurs de $n$ tandis que
$f(n/2)$ peut y être de l'ordre de $(n/2)^{2}$ : l'inégalité
$2f(n/2)\le c\,f(n)$ est alors violée, le coût d'un niveau inférieur dépasse celui de la
racine, et la conclusion $T(n)=\Theta(f(n))$ tombe en défaut. La condition de régularité
est ce qui interdit ces oscillations pathologiques (elle est automatiquement vérifiée
pour $f(n)=n^{\beta}\log^{\gamma}n$, cas de toutes les récurrences rencontrées en
pratique).

% ---------------------------------------------------------------------
\exo{Multiplications rapides : gagner un appel récursif}
\reponse
\textbf{1. Karatsuba.}
\begin{enumerate}[label=(\alph*),leftmargin=*,itemsep=4pt]
  \item Avec $x=x_1 10^{m}+x_0$ et $y=y_1 10^{m}+y_0$ ($m=n/2$) :
        \[
          xy = x_1y_1\,10^{2m} + (x_1y_0+x_0y_1)\,10^{m} + x_0y_0 .
        \]
        L'algorithme naïf calcule les $4$ produits $x_1y_1$, $x_1y_0$, $x_0y_1$,
        $x_0y_0$ de nombres à $n/2$ chiffres, puis effectue additions et décalages en
        $\Theta(n)$ :
        \[
          T(n)=4\,T(n/2)+\Theta(n).
        \]
        Ici $\alpha=\log_2 4 = 2$ et $f(n)=n=O(n^{2-\varepsilon})$ : cas 1, donc
        $T(n)=\Theta(n^{2})$ --- on ne gagne rien sur la multiplication posée.
  \item $(x_1+x_0)(y_1+y_0) = x_1y_1 + x_1y_0 + x_0y_1 + x_0y_0$, d'où
        \[
          x_1y_0+x_0y_1 = (x_1+x_0)(y_1+y_0) - x_1y_1 - x_0y_0 .
        \]
        Les trois produits $P_1=x_1y_1$, $P_2=x_0y_0$ et $P_3=(x_1+x_0)(y_1+y_0)$
        suffisent : $xy = P_1 10^{2m} + (P_3-P_1-P_2)10^{m}+P_2$.
        \textbf{Trois} multiplications récursives au lieu de quatre (les additions
        supplémentaires restent en $\Theta(n)$).
  \item $T(n)=3\,T(n/2)+\Theta(n)$ : $\alpha=\log_2 3\approx 1{,}585$ et
        $f(n)=n=O(n^{\alpha-\varepsilon})$ : cas 1, donc
        \[
          \boxed{T(n)=\Theta\!\left(n^{\log_2 3}\right)=\Theta\!\left(n^{1{,}585}\right)},
          \qquad \log_2 3 = 1{,}585 \text{ à } 10^{-3} \text{ près}.
        \]
        Pour $n=10^{6}$ chiffres, le rapport $n^{2}/n^{1,585}=n^{0,415}\approx 300$ :
        le gain est considérable.
\end{enumerate}

\textbf{2. Strassen.}
\begin{enumerate}[label=(\alph*),leftmargin=*,itemsep=4pt]
  \item Le produit par blocs naïf effectue $8$ produits de blocs $n/2\times n/2$ et
        $\Theta(n^2)$ additions : $T(n)=8T(n/2)+\Theta(n^{2})$.
        $\alpha=\log_2 8 = 3$, $f(n)=n^{2}=O(n^{3-\varepsilon})$ : cas 1, donc
        $T(n)=\Theta(n^{3})$. On retrouve bien la complexité de l'algorithme usuel
        à trois boucles imbriquées.
  \item Strassen : $T(n)=7T(n/2)+\Theta(n^{2})$, $\alpha=\log_2 7\approx2{,}807$,
        $f(n)=n^2=O(n^{\alpha-\varepsilon})$ : cas 1, donc
        \[
          \boxed{T(n)=\Theta\!\left(n^{\log_2 7}\right)=\Theta\!\left(n^{2{,}807}\right)}.
        \]
        Gagner un seul produit sur huit fait passer l'exposant de $3$ à $2{,}807$.
  \item Pour $T(n)=a\,T(n/2)+\Theta(n^{2})$ on a $\alpha=\log_2 a$ et l'on compare
        $\alpha$ à $2$ :
        \begin{itemize}[leftmargin=*,itemsep=1pt]
          \item $a<4$ : $\alpha<2$, cas 3, $T(n)=\Theta(n^{2})$ ;
          \item $a=4$ : $\alpha=2$, cas 2, $T(n)=\Theta(n^{2}\log n)$ ;
          \item $a>4$ : $\alpha>2$, cas 1, $T(n)=\Theta(n^{\log_2 a})$.
        \end{itemize}
        La récurrence cesse donc d'être en $\Theta(n^{2}\log n)$ dès que $a\ne4$, et
        c'est à partir de $a=5$ que la complexité dépasse strictement
        $\Theta(n^{2}\log n)$.
\end{enumerate}

% ---------------------------------------------------------------------
\exo{Changement de variable}
\reponse
\textbf{1.} On pose $n=2^{k}$ et $S(k)=T(2^{k})$. Alors $\sqrt n = 2^{k/2}$ et
\[
  S(k)=S(k/2)+1,\qquad S(1)=T(2)=1 .
\]
C'est la récurrence de la recherche dichotomique : $S(k)=\log_2 k+1$.
Comme $k=\log_2 n$ :
\[
  T(n) = \log_2\log_2 n + 1 = \Theta(\log\log n).
\]

\textbf{2.} Même changement de variable : $\log_2 n = k$, donc
\[
  S(k)=2\,S(k/2)+k,\qquad S(1)=1 .
\]
Théorème général en la variable $k$ : $a=2$, $b=2$, $\alpha=1$, $f(k)=k=\Theta(k^{1})$ :
cas 2 avec $c=0$, donc $S(k)=\Theta(k\log k)$. En revenant à $n$ :
\[
  \boxed{T(n)=\Theta\!\big(\log n \cdot \log\log n\big)} .
\]

\textbf{3.} $T(n)=\sqrt n\,T(\sqrt n)+n$ devient, avec $S(k)=T(2^{k})$ :
\[
  S(k) = 2^{k/2}\,S(k/2) + 2^{k}.
\]
On divise par $2^{k}$ et l'on pose $U(k)=S(k)/2^{k}$ :
\[
  U(k) = \frac{2^{k/2}S(k/2)}{2^{k}} + 1 = \frac{S(k/2)}{2^{k/2}} + 1 = U(k/2)+1 .
\]
Donc $U(k)=\Theta(\log k)$, puis $S(k)=\Theta\!\big(2^{k}\log k\big)$ et enfin
\[
  \boxed{T(n)=\Theta\!\big(n\log\log n\big)} .
\]

\textbf{4.} La récurrence est $T(n)=T(\sqrt n)+\Theta(n)$. Avec $S(k)=T(2^{k})$ :
$S(k)=S(k/2)+\Theta(2^{k})$. En déroulant, le coût total est
\[
  \Theta\!\left(2^{k}+2^{k/2}+2^{k/4}+\cdots\right)
  = \Theta\!\left(2^{k}\right),
\]
car la somme est dominée par son premier terme (le second est déjà $\sqrt{2^k}$).
Donc $T(n)=\Theta(n)$ : le coût est entièrement concentré dans le \emph{premier}
niveau de recombinaison. C'est l'exact opposé du cas 1 du théorème général, où le coût
se concentre dans les feuilles.

\vspace{14pt}
{\color{teal}\hrule height 1pt}
\vspace{6pt}
\begin{center}
\small\itshape Fin du corrigé --- 16 exercices.
\end{center}

\end{document}
